\setcounter{chapter}{4}\chapter{Fonctions trigonométriques}\label{ch:5}

{\small\sffamily\color{teal}S75, S76, S78, S79\par}

\begin{retenir}{Objectifs}Utiliser parité et périodicité ; dériver sinus et cosinus ; résoudre équations et inéquations ; étudier une fonction périodique.\end{retenir}

\section{Notes de cours}
\subsection{Cercle et fonctions}
Sur le cercle trigonométrique, le point repéré par un angle $x$ en radians a pour coordonnées $(\cos x,\sin x)$. Pour tout réel $x$,
\[\cos^2x+\sin^2x=1,\quad -1\le\cos x,\sin x\le1.\]
Cosinus est paire : $\cos(-x)=\cos x$. Sinus est impaire : $\sin(-x)=-\sin x$. Les deux fonctions sont $2\pi$-périodiques.
\[\begin{array}{c|ccccc}x&0&\frac{\pi}{6}&\frac{\pi}{4}&\frac{\pi}{3}&\frac{\pi}{2}\\\hline
\cos x&1&\frac{\sqrt{3}}{2}&\frac{\sqrt{2}}{2}&\frac{1}{2}&0\\
\sin x&0&\frac{1}{2}&\frac{\sqrt{2}}{2}&\frac{\sqrt{3}}{2}&1
\end{array}\]
\subsection{Dérivées et variations}
\[(\sin x)'=\cos x,\qquad(\cos x)'=-\sin x.\]\[
(\sin(ax+b))'=a\cos(ax+b),\quad(\cos(ax+b))'=-a\sin(ax+b).\]
Sur $[-\pi,\pi]$, cosinus croît sur $[-\pi,0]$ de $-1$ à $1$, puis décroît sur $[0,\pi]$ de $1$ à $-1$. Sinus décroît sur $[-\pi,-\frac{\pi}{2}]$, croît sur $[-\frac{\pi}{2},\frac{\pi}{2}]$, puis décroît sur $[\frac{\pi}{2},\pi]$.
\paragraph{Justification.} Lire le signe de sinus et cosinus sur le cercle puis appliquer le lien dérivée--variations. La périodicité prolonge les courbes ; elle ne dispense pas d'étudier un intervalle de longueur suffisante.
\subsection{Équations et inéquations}
Si $|a|>1$, $\cos x=a$ et $\sin x=a$ n'ont aucune solution réelle. Pour un angle connu $\theta$,
\[\cos x=\cos\theta\Longleftrightarrow x=\theta+2k\pi\ \text{ou}\ x=-\theta+2k\pi,\]
\[\sin x=\sin\theta\Longleftrightarrow x=\theta+2k\pi\ \text{ou}\ x=\pi-\theta+2k\pi,\quad k\in\Z.\]
On restreint ensuite les solutions à l'intervalle demandé et on élimine les doublons.
\paragraph{Exemples sur $[-\pi,\pi]$.} $\cos x=\frac{1}{2}$ donne $x=\pm\frac{\pi}{3}$ ; $\sin x=\frac{1}{2}$ donne $x=\frac{\pi}{6}$ ou $\frac{5\pi}{6}$. $\cos x\le\frac{1}{2}$ donne $[-\pi,-\frac{\pi}{3}]\cup[\frac{\pi}{3},\pi]$ ; $\sin x\ge\frac{1}{2}$ donne $[\frac{\pi}{6},\frac{5\pi}{6}]$.
\begin{methode}{Une composée trigonométrique}
Pour résoudre $\cos (2x)=\frac{1}{2}$ sur $[-\pi,\pi]$, poser $t=2x\in[-2\pi,2\pi]$, résoudre pour $t$, puis diviser par $2$. On obtient $x\in\{-\frac{5\pi}{6},-\frac{\pi}{6},\frac{\pi}{6},\frac{5\pi}{6}\}$.
\end{methode}
\paragraph{Optimisation.} Pour $f(x)=\sin x+\frac12\sin(2x)$ sur $[0,\pi]$, $f'=\cos x+\cos(2x)=2\cos^2x+\cos x-1=(2\cos x-1)(\cos x+1)$. La fonction croît jusqu'à $\frac{\pi}{3}$ puis décroît ; maximum $\frac{3\sqrt{3}}{4}$.
\paragraph{Identités de duplication utilisées dans les exercices.} Pour tout réel $x$, $\sin(2x)=2\sin x\cos x$ et $\cos(2x)=\cos^2x-\sin^2x=2\cos^2x-1$. Ces identités sont rappelées ici pour que leur emploi dans les problèmes ne soit pas implicite.
\paragraph{Symétries des courbes.} La parité de cosinus correspond à une symétrie par rapport à l'axe des ordonnées ; l'imparité de sinus à une symétrie centrale de centre l'origine. La périodicité prolonge chaque courbe par translation de vecteur $(2\pi,0)$.


\begin{center}\begin{tikzpicture}[scale=1.6]
\draw[->,gray](-1.3,0)--(1.4,0);\draw[->,gray](0,-1.25)--(0,1.4);\draw[navy](0,0)circle(1);
\draw[teal,thick](0,0)--(.5,.866);\draw[teal,dashed](.5,0)--(.5,.866)--(0,.866);
\draw[->,teal](.35,0)arc(0:60:.35);\node[right]at(.25,.27){$x$};
\node[below]at(.5,0){$\cos x$};\node[left]at(0,.866){$\sin x$};\node[above right]at(.5,.866){$M$};
\end{tikzpicture}\par\small Les coordonnées du point du cercle sont $(\cos x,\sin x)$.\end{center}

\section{Activité de recherche}\textbf{Hauteur d'un point sur une roue --- 35 min.}\par Une roue tourne à vitesse constante dans un plan vertical. Son rayon est de $2$ m et son centre se trouve à $3$ m du sol. Un point de la roue est à sa position la plus basse à l'instant $t=0$. La roue effectue un tour en $8$ s. Sa hauteur est modélisée par $h(t)=3-2\cos\left(\frac{\pi t}{4}\right)$, où $t$ est exprimé en secondes et $h(t)$ en mètres.
\begin{enumerate}\item Vérifier que le modèle donne la hauteur initiale annoncée, une période de $8$ s et les hauteurs minimale et maximale attendues.
\item Calculer $h' (t)$, puis étudier les variations de $h$ sur $[0;8]$.
\item Déterminer tous les instants $t\in[0;8]$ pour lesquels la hauteur vaut $4$ m.
\item Déterminer la durée totale, au cours de ce tour, pendant laquelle le point se trouve à au moins $4$ m du sol.\end{enumerate}

\section{Exercices progressifs}

\exercice{1}{Entraînement}Résoudre dans $[-\pi;\pi]$ l'équation $\sin x=-\frac12$.

\exercice{2}{Entraînement}Résoudre dans $[-\pi;\pi]$ l'inéquation $\cos x\ge\frac{\sqrt2}{2}$.

\exercice{3}{Entraînement}On définit sur $\R$ les fonctions $f(x)=\cos(3x-1)$ et $g(x)=\sin^2x$. Calculer leurs dérivées.

\exercice{4}{Entraînement}On considère $f(x)=2+\sin x$ sur $[0;2\pi]$. Calculer sa dérivée et dresser son tableau de variations.

\exercice{5}{Entraînement}Calculer $\displaystyle\int_0^\pi\sin x\,\dd x$, puis la valeur moyenne de la fonction sinus sur $[0;\pi]$.

\section{Travail dans le manuel ou une annale}\textbf{Barbazo, p. 166, exercice 19.} En S75 : question 1, pour lire une symétrie et une période sur une courbe. À la maison : question 2, pour démontrer les conjectures par des identités valables pour tout réel.

\section{Devoir à la maison}\textbf{DM5 --- Partager une aire --- 55 min, après S79.}\par On considère la fonction $f$ définie sur $[0;\frac\pi2]$ par $f(x)=\sin x\cos x$. Sa courbe est représentée dans un repère orthonormé.
\begin{enumerate}\item Calculer $f' (x)$ et déterminer le maximum de $f$ sur cet intervalle.
\item Résoudre sur $[0;\frac\pi2]$ l'inéquation $f(x)\ge\frac14$.
\item Calculer $\displaystyle\int_0^{\frac{\pi}{2}}f(x)\,\dd x$, puis la valeur moyenne de $f$ sur $[0;\frac\pi2]$.
\item Pour $a\in[0;\frac\pi2]$, on pose $A(a)=\displaystyle\int_0^a f(x)\,\dd x$. Déterminer $a$ pour que le domaine sous la courbe entre $0$ et $a$ ait une aire égale à la moitié de l'aire totale sous la courbe entre $0$ et $\frac\pi2$.\end{enumerate}

\ifprof\clearpage\section{Conduite et corrigés}\begin{profbox}{Repères pour le professeur}Ce chapitre est réellement terminé début mai, et non fin avril. Il n'entre donc pas dans le deuxième blanc d'avril. Reprendre le cercle avant les courbes. Utiliser les radians dans les dérivées. Faire contrôler les bornes $-\pi$ et $\pi$ : elles représentent le même point du cercle mais sont deux nombres distincts si l'intervalle les contient. Les variations de composées se déduisent du signe complet de la dérivée.\end{profbox}

\subsection{Corrigé de l'activité}\begin{enumerate}
\item $h(0)=3-2=1$ m. Ajouter $8$ à $t$ ajoute $2\pi$ à l'argument du cosinus : $8$ s est une période, et c'est la plus petite période positive. Comme $-1\le\cos\le1$, $1\le h\le5$ ; les deux bornes sont atteintes.
\item Par dérivation d'une composée, $h'(t)=(\frac{\pi}{2})\sin (\frac{\pi t}{4})$. L'argument parcourt $[0,2\pi]$. Le sinus est positif sur $]0,\pi[$, puis négatif sur $]\pi,2\pi[$. Ainsi $h$ croît de $1$ à $5$ m sur $[0,4]$, puis décroît de $5$ à $1$ m sur $[4,8]$.
\item $h(t)=4\Longleftrightarrow\cos (\frac{\pi t}{4})=-\frac{1}{2}$. Sur $[0,2\pi]$, les angles correspondants sont $\frac{2\pi}{3}$ et $\frac{4\pi}{3}$ ; donc $t=\frac{8}{3}$ s ou $t=\frac{16}{3}$ s.
\item La condition $h(t)\ge4$ équivaut à $\cos (\frac{\pi t}{4})\le-\frac{1}{2}$, donc à $t\in[\frac{8}{3},\frac{16}{3}]$. La durée est $\frac{16}{3}-\frac{8}{3}=\frac{8}{3}$ s.
\end{enumerate}

\paragraph{E1.}Sur le cercle, le sinus vaut $-\frac{1}{2}$ aux angles $-\frac{\pi}{6}$ et $-\frac{5\pi}{6}$ dans $[-\pi,\pi]$. Les autres angles s'en déduisent par ajout de $2k\pi$ et sont hors de l'intervalle demandé.

\paragraph{E2.}Le cosinus est au moins $\frac{\sqrt{2}}{2}$ sur l'arc autour de l'angle $0$, entre $-\frac{\pi}{4}$ et $\frac{\pi}{4}$, bornes incluses. L'ensemble des solutions est donc $[-\frac{\pi}{4},\frac{\pi}{4}]$.

\paragraph{E3.}Par composition, $(\cos(3x-1))'=-3\sin(3x-1)$. Pour le carré du sinus, $(\sin^2x)'=2\sin x\cos x$. Les deux dérivées sont définies sur $\R$.

\paragraph{E4.}$f'(x)=\cos x$, positive sur $[0,\frac{\pi}{2}[$ et $]\frac{3\pi}{2},2\pi]$, négative sur $]\frac{\pi}{2},\frac{3\pi}{2}[$. Ainsi $f$ croît de $2$ à $3$, décroît de $3$ à $1$, puis croît de $1$ à $2$. Le maximum est $3$ en $\frac{\pi}{2}$ et le minimum $1$ en $\frac{3\pi}{2}$.

\paragraph{E5.}Une primitive de sinus est $-\cos x$. Donc l'intégrale vaut $[-\cos x]_0^\pi=1-(-1)=2$. La valeur moyenne sur cet intervalle de longueur $\pi$ est $\frac{2}{\pi}$.

\subsection{Corrigé du DM5}\begin{enumerate}
\item Par produit, $f'=\cos^2x-\sin^2x=\cos2x$. Sur l'intervalle étudié, $2x\in[0,\pi]$ : la dérivée est positive avant $x=\frac{\pi}{4}$ et négative après. Le maximum vaut $f(\frac{\pi}{4})=(\frac{\sqrt{2}}{2})^2=\frac{1}{2}$.
\item Grâce à $\sin2x=2\sin x\cos x$, $f(x)\ge\frac{1}{4}$ équivaut à $\sin 2x\ge\frac{1}{2}$. Pour $2x\in[0,\pi]$, cela donne $2x\in[\frac{\pi}{6},\frac{5\pi}{6}]$, donc $x\in[\frac{\pi}{12},\frac{5\pi}{12}]$.
\item La dérivée de $\frac12\sin^2x$ est $\sin x\cos x=f(x)$. Donc $\int_0^{\frac{\pi}{2}}f(x)\,\dd x=[\frac12\sin^2x]_0^{\frac{\pi}{2}}=\frac{1}{2}$. La valeur moyenne est $\frac{\frac{1}{2}}{\frac{\pi}{2}}=\frac{1}{\pi}$.
\item $A(a)=\frac12\sin^2a$. La moitié de l'aire totale est $\frac{1}{4}$, donc $\sin^2a=\frac{1}{2}$. Sur $[0,\frac{\pi}{2}]$, sinus est positif et strictement croissant : cette égalité équivaut à $\sin a=\frac{\sqrt{2}}{2}$, puis à $a=\frac{\pi}{4}$.
\end{enumerate}

\subsection{Corrigé du travail dans le manuel ou une annale}La courbe de $f(x)=\cos(3x)$ suggère une symétrie par rapport à l'axe des ordonnées et une répétition après une translation horizontale de $\frac{2\pi}3$. Pour tout $x\in\R$, $f(-x)=\cos(-3x)=\cos(3x)=f(x)$ : $f$ est paire. D'autre part,
\[f\left(x+\frac{2\pi}3\right)=\cos(3x+2\pi)=\cos(3x)=f(x).\]
Le nombre $\frac{2\pi}3>0$ est donc une période de $f$. Le dessin suggère ces propriétés ; les identités les établissent pour tout réel.

\fi

\backmatter
